% SCRATCH DRAFT — Vol 2 courtroom coda, Ch25 "Resolution, Precision, and the Continuum" (task vol2-coda, author: Podo).
% THIS IS NOT THE BOOK. Do not \input. FLOOR BANKED (integrate on Kodo's accept). Fresh-append (no old beat; body ends line 168).
%
% ⚠ FENCES:
%  - The continuum = the SMOOTH-SHADOW analogy, OPEN since Ch1: a shadow the discrete casts as it refines, never
%    the substrate/ground; a limit approached, not a floor stood upon; the most useful analogy + the most fenced.
%    Work this fully. BUT ⚠ RESERVE the deepest punch for Ch30: do NOT say "the continuity of time rests on
%    nothing", "never observed / never derived", "named/adopted as an assumption", or that the physics secretly
%    rests on it. Coda stays at "shadow not ground"; the "unfounded assumption the edifice stands on" is Ch30's.
%  - The irrational (√2, the plain square's diagonal) = the chapter's limit: seed "the ledger can approach but
%    never finish, must take a smallest step" — SEED ONLY; do NOT open forcing / Cantor-diagonal-of-the-reals /
%    metamath (Ch29). box/type-identity RESERVED. ±1/spinor OUT.
%  - Ch26 handoff (determinism/induction/cause): ⚠ free-will WITHHELD — seed "whether the past fixes the future /
%    the gap between past and future", NOT choice/agency.
%  - Charge re-sounds (count read to a resolution floor; continuum a shadow of the count). Consensus consolidate
%    (agreement reserved Ch30). "verdict"/2nd-person/penalty-"fine"/free-will ABSENT.
%
% CONTINUITY: state carries from Ch24. Ch24 closed: "how sharp a measurement can be made ... what the sharpest
% reading of all still cannot reach: the continuum, and the limits of reading it." Ch25 IS that continuum.
%
% PROPOSED SEAM in books/experimentation/latex/chapters/25.tex:
%   KEEP the entire body through line 168 ("...and what the gap between past and future costs.").
%   APPEND: \bigskip / \begin{center}\rule{0.4\textwidth}{0.4pt}\end{center} / \bigskip / then the coda below. (Fresh append.)
%
% >>>>>>>>>>>>>>>>>>>>>> APPEND BEGINS (fresh \rule + coda after the body) >>>>>>>>>>>>>>>>>>>>>>

\bigskip
\begin{center}
\rule{0.4\textwidth}{0.4pt}
\end{center}
\bigskip

\noindent The case asks, before the end, how finely its readings can be made --- how sharply the world can be told
apart --- and in the asking it meets the one thing it has held at arm's length the whole book long: the
continuum, the unbroken line physics is written on. It is the instrument's most useful companion and the one it
has handled most warily, for the continuum is never a thing the instrument holds; it is a shadow the instrument
casts.

Every reading has a floor. Below some smallest step the instrument records nothing --- two states nearer than
that are, to it, one state, and no care sharpens what lies beneath the floor; a finer ruler only lowers the
floor, it does not abolish it. This is the resolution the very first chapter laid down: the grid we set, below
which nothing varies. Galileo, grinding lenses to push the floor down, saw moons around Jupiter the naked eye
could not --- not because the moons were tiny beyond measure but because they lay below the eye's floor and
above the glass's; and Berkeley's mockery of the infinitesimal, the "ghosts of departed quantities," is answered
by exactly this floor: there are no departed quantities, only distinctions too small to count. And what the instrument returns above the floor is precise when it returns the same shadow
each time --- but precision is not truth: a tight cluster of readings can sit, repeatably, off the mark, and
whether they are right is a separate question the sameness never answers. The smooth curve the instrument draws
through its readings is just that --- a curve drawn \emph{through}, the shape the discrete anchors approach, and
not a thing lying beneath them. Even how long a rough thing is depends on the ruler that measures it: shrink the
yardstick and a coastline's length climbs without end, for there was never one length, only a length at a scale
--- how long is the coast of Britain, the famous question runs, and the honest answer is that there is no
resolution-free length to give.

So the continuum stands where it has stood since the first count: the smooth shadow the discrete casts as it is
refined, never the ground the discrete rests on. It is the most useful analogy the instrument keeps --- the
bridge to every equation, drawn again and again --- and the most carefully fenced, offered always as a
convenience of ours and never as a fact the world hands over. Press the ledger to write out an irrational ---
the diagonal of a plain square, which no ratio of whole numbers can ever equal --- and the refinement never
finishes: the ledger cannot reach it, only approach, and must at the last take a smallest step and stop. Legend says the
first man to prove that diagonal beyond the reach of any ratio drowned for telling of it, so hard did the old
creed that all things are whole numbers die. The
instrument names the continuum, uses it, leans its whole calculus on it --- and marks, each time, that it is a
shadow named as a shadow, a limit approached and not a floor to stand on. What more there is to say of it, the
case keeps for the end.

And the \emph{charge}, read this finely, meets a floor of its own. The count the case will collect on is read to
a smallest step and no further; below it two readings are one, and no sharpening tells them apart. The charge is
exact to the floor and silent beneath it --- a count in whole cells, not a real number carried out forever. It
is never owed to more places than the reading has; a charge claimed sharper than the ruler that took it would be
a precision the record does not hold, and the instrument will not pretend to it. And
the continuum that would smooth the count into an unbroken quantity is, here as everywhere, a shadow of the
count and not its ground: the tally is the real thing, the smooth curve through it a convenience. The trace,
likewise, is a count kept to its floor --- what the record retains, retained in whole steps, exact as far as it
reads and honest about where it stops. Whatever comes to be owed is owed in counted cells, to the floor of the
reading and no finer.

And each of the three readings the case will bring has a floor and a precision of its own --- its own smallest
step, its own reproducibility. To meet on one number the three must agree not to a perfection none of them holds
but to the coarsest floor among them; agreement finer than the roughest ruler would be a fiction. The three are
corrected and ready, each honest about how finely it reads; the number they must meet on is still not in hand. The
floor of the reading is one more thing the meeting will have to respect.

The instrument has one inward question left before the reckoning. Having asked how finely a record can be read, it
asks next whether a record can be carried forward at all --- whether a finite past fixes its future or leaves a
gap between them, and what that gap, if it is there, costs. The next chapter turns to determinism, induction,
and cause. The case is still open, its readings admissible and floored and not yet agreed; the court does not
sit until the record is complete.

% <<<<<<<<<<<<<<<<<<<<<< APPEND ENDS <<<<<<<<<<<<<<<<<<<<<<
